% Avsnitt 18.11, ur lectures/L18/appendix/b_exercises.md.

\section{Övningar}
\label{c18:sec:exercises}

\begin{aside}
Skriv din egen \file{register_bank.vhd} utifrån \secref{c18:sec:iface} till \ref{c18:sec:write}; den
utdelade testbänken, körd enligt bilaga~\ref{app:sim}, är kontrollen av att den beter sig korrekt.
Övning~\ref{c18:ex:waveform} och \ref{c18:ex:slave} handlar om \code{spi_slave.vhd}, som delas ut i
stället för att skrivas \textemdash{} den läser du, och \chapref{c20:ch} förväntar sig att du kan
den.

Den auktoritativa källan för varje registernummer nedan är bilaga~\ref{app:protocol}. Där en övning
och en specifikation är oense vinner specifikationen, och det är övningen som är buggen.
\end{aside}

\exercise{\code{register_bank}}{VHDL}
\label{c18:ex:bank}
Skriv din egen \file{register_bank.vhd} utifrån kapitlet. Verifiera den:

\begin{shell}
cd bridge
ghdl -a --std=93 ../controller/can_def.vhd register_bank.vhd register_bank_tb.vhd
ghdl -e --std=93 register_bank_tb
ghdl -r --std=93 register_bank_tb --assert-level=error
\end{shell}

Om en kontroll fallerar, läs assertion-meddelandet; vart och ett av de elva fallen säger vad ett
haveri i just det fallet betyder.

\exercise{De tre STATUS-bitarna}{Resonemang}
\label{c18:ex:status}
\xpart{a} Ur reset läser STATUS \code{0x00000001}. Säg vilken bit det är och varför resetvärdet är
\code{'1'} och inte \code{'0'}. Beskriv sedan vad en drivrutins \code{send()} skulle göra, i all
evighet, om det vore \code{'0'}.

\xpart{b} Bit 2 låser en \textbf{stigande flank} hos kontrollerns \code{error}, inte dess nivå. Anta
att den i stället låste nivån. Skriv ner den följd av drivrutinsoperationer \textemdash{} polla,
nollställ, polla \textemdash{} som då skulle bete sig annorlunda, och säg vad drivrutinens författare
skulle dra för slutsats om vad som var trasigt.

\xpart{c} Bit 1 sätts av \code{rx_valid} och nollställs av en skrivning till \code{RX_ACK}. Bit 0:s
huvudsakliga sättvillkor är \code{tx_done}, och den nollställs av en skrivning till \code{TX_SEND}.
Det ena av de paren har sina villkor för att sätta och nollställa i den ”naturliga” ordningen, och
det andra är omvänt. Säg vilket som är vilket, och varför det omvända ändå är rätt.

\xpart{d} \Secref{c18:sec:status} ger en ordningsregel: kontrollerhändelser tillämpas \textbf{efter}
bryggans skrivningar i den klockade processen. Konstruera den enda klockflank där den regeln ändrar
utfallet, namnge båda händelserna, och säg vilken av de två möjliga lösningarna som kan låsa systemet
och vilken som inte kan det.

\xpart{e} Bit 0 har \emph{tre} sättvillkor, och ett av dem är en stigande flank hos \code{error} medan
bit 0 redan är låg. Förklara varför det villkoret inte behöver veta något om huruvida kontrollern
sände eller tog emot \textemdash{} vad är det i en nods beteende som gör att ”error steg medan vi var
upptagna” betyder ”en sändning tog just slut”?

\xpart{f} Ta bort det villkoret ur din bank och kör om testbänken. Vilket fall fallerar, och vilken
konsekvens säger dess meddelande att det skulle ha fått? Förklara sedan varför det här är skydd i
flera lager snarare än dubbelarbete, givet att \chapref{c17:ch} redan låter kontrollern pulsa
\code{tx_done} vid en avbruten sändning. (Vem skriver \code{can_controller}? Vem skrev testbänken som
skulle fånga det?)

\xpart{g} \code{TX_ABORT} är den tredje vägen, och till skillnad från de andra två motsvarar den inte
något hårdvaran gjorde \textemdash{} det är drivrutinen som går förbi banken. Säg vad den \emph{inte}
gör: namnge minst två saker du kanske skulle vänta dig att ett register som heter \code{TX_ABORT} når
fram till, och förklara varför den inte når någondera. (Ledtråden är \code{can_controller}s portlista,
som fastställdes i \chapref{c11:ch}.)

\exercise{\code{TX_SEND} och pulsdisciplinen}{Simulering}
\label{c18:ex:txsend}
\xpart{a} En skrivning av \code{0x1} till \code{TX_SEND} höjer \code{tx_req} under exakt en cykel.
Säg exakt vilken cykel, i förhållande till den klockflank som bar \code{reg_write}, och vilka två
saker \code{register_bank_tb} kontrollerar om den.

\xpart{b} Gör nu sönder det med flit. Ändra din \code{TX_SEND}-gren så att \code{tx_req} blir en
\textbf{nivå}: sätt den vid en accepterad skrivning och nollställ den först när \code{tx_done}
kommer. Kör om \code{register_bank_tb}. Vilken kontroll fallerar först, och vad säger dess meddelande?

\xpart{c} Din trasiga version från \textbf{b)} skulle fortfarande \emph{se} korrekt ut på ett
oscilloskop som tittar på en enda ram. Beskriv, med \secref{c16:sec:roles}:s regel för när
\code{can_controller} agerar på \code{tx_req}, vad den gör vid den andra ramen och varje ram
därefter. Säg sedan varför det här är en bugg som simulering fångar lätt och ett test på bänken inte
gör.

\xpart{d} En skrivning till \code{TX_SEND} medan STATUS bit 0 är låg kastas helt \textemdash{} ingen
puls, ingen tillståndsändring. En annan design skulle i stället kunna köa den och avfyra när
kontrollern rapporterar redo. Ge ett konkret argument för regeln att kasta den, givet att drivrutinen
kan polla STATUS. Ställ tillbaka modulen som den var.

\exercise{Att läsa av en SPI-transaktion ur vågformen}{Resonemang}
\label{c18:ex:waveform}
Det här kapitlet heter ”från vågformen” av ett skäl: \chapref{c20:ch} kan mycket väl räcka dig en och
fråga vad den betyder. Läge 0 rakt igenom \textemdash{} SCK vilar lågt, MOSI samplas på
\textbf{stigande} flank, MISO ändras på \textbf{fallande} flank, MSB först.

\xpart{a} Rita för hand hela fembytestransaktionen som skriver \code{0x00000001} till
\code{TX_SEND}. Räkna ut kommandobyten ur formatet i bilaga~\ref{app:protocol} i stället för att
kopiera den: bit 7 är skrivbiten, bit 6--4 är reserverade, bit 3--0 är registerindexet, och
\code{TX_SEND} är index 5. Märk upp \code{SS}, \code{SCK} och \code{MOSI}, och markera de 40 stigande
flankerna i grupper om åtta.

\xpart{b} Markera på samma ritning var mastern inte får släppa \code{SS}, och vad slaven måste göra
om den ändå gör det. Återge regeln ur protokollspecifikationens avsnitt om avbrott med egna ord.

\xpart{c} En läsning av \code{STATUS} är spegelbilden: kommandobyte \code{0x00}, fyra dummybytes ut,
fyra bytes registervärde tillbaka på MISO. Protokollspecifikationen säger att slaven låser
registervärdet \textbf{en gång}, vid slutet av kommandobyten. Förklara vad en drivrutin skulle kunna
observera om slaven i stället läste registret på nytt för var och en av de fyra bytesen, och varför
just \code{STATUS} är det register där det skulle göra ont.

\xpart{d} Räkna SCK-flankerna i en hel transaktion och räkna, vid det maximum på 1~MHz som
protokollspecifikationen tillåter, ut hur lång tid en transaktion tar. Jämför det med den 20~ns långa
puls som registerbanken finns till för att översätta. Det förhållandet är hela argumentet för det här
lagret, i ett enda tal.

\exercise{Att läsa \code{spi_slave.vhd}}{Resonemang}
\label{c18:ex:slave}
Du skriver inte den här modulen, men du måste kunna läsa den.

\xpart{a} \code{spi_slave} använder inte \code{SCK} som klocka. Den synkroniserar in den i
50 MHz-domänen och detekterar dess flanker i stället. Ge två skilda skäl, ett om metastabilitet
(\chapref{c4:ch}) och ett om vad en design med två orelaterade klockor kostar i en FPGA.

\xpart{b} Posten \code{sync_t} har tre fält, men bara två av dem är synkroniseringssteg. Säg vad det
tredje är till för, och varför \code{mosi} behöver två fält där \code{sclk} och \code{ss} behöver
tre.

\xpart{c} I resetgrenen nollställs \code{sclk_s} till idel nollor och \code{ss_s} till idel
\textbf{ettor}. Förklara varför de skiljer sig, och följ sedan vad \code{ss_active} skulle läsa under
de två första klockflankerna efter reset om \code{ss_s} nollställdes till nollor utan någon master
inkopplad. Det här är samma regel som Övning~\ref{c4:ex:buttonsync} byggdes kring; säg vilken kontroll
i den övningens testbänk som motsvarar den.

\xpart{d} Hanteringen av SCK-flanker ligger inuti \code{else}-armen i ett test på \code{ss_s.s2}, så
en avvald slav ignorerar SPI-ledningarna helt. Konstruera felet som skulle följa om den inte gjorde
det: utgå från en vilsen SCK-flank medan \code{SS} är hög, och följ den fram till vad \emph{nästa}
registerskrivning hamnar på.

\xpart{e} Grenen \code{ss_falling} är ett \code{elsif} före SCK-grenarna i stället för ett separat
\code{if} bredvid dem. Säg vid vilken enda klockflank den ordningen spelar roll, och vad som skulle
vara fel med de mottagna bytesen om de två grenarna kördes i omvänd ordning.

\exercise{Infångningen, och gapet den inte täpper till}{Resonemang}
\label{c18:ex:capture}
\xpart{a} När \code{rx_valid} pulsar fångar banken in \code{rx_id}, \code{rx_dlc} och \code{rx_data}
i \code{RX_*}-registren i stället för att koppla kontrollerns portar rakt in i läsmultiplexern.
Förklara vad en drivrutin som läser fyra separata SPI-transaktioner skulle se utan den infångningen.

\xpart{b} En skrivning till \code{RX_ACK} nollställer STATUS bit 1 men lämnar \code{RX_*}-registren
med sitt innehåll kvar. Säg varför det är ofarligt, och citera meningen i registerkartan som gör det
så.

\xpart{c} En andra ram anländer innan den första har kvitterats. Säg exakt vad drivrutinen
observerar, och varför det här är ett dokumenterat gap som ärvts från \code{can_controller}
(\chapref{c11:ch}) snarare än en brist hos registerbanken.

\xpart{d} Konstruera, enbart på papper, den minsta ändring i \emph{registerkartan} som skulle låta en
drivrutin upptäcka att den missat en ram. Du lägger inte till någon mottagningskö; du lägger till
möjligheten att veta att en behövdes. Säg vad det kostar i hårdvara, och i vilket befintligt register
det skulle kunna bo utan att bryta det kontrakt parallellklassen redan skriver mot.
